\ifdefined\iscombined
	\lecturetitle{Transfer Functions from Circuits}
\else
\documentclass{article}
\usepackage{xparse}
\usepackage{changepage}
\usepackage{listings}
\usepackage{amsthm}
\usepackage{comment}
\usepackage{amsmath}
\usepackage{braket}
\usepackage{amsfonts}
\usepackage{sectsty}
\usepackage{hyperref}
\usepackage{fancyhdr}
\usepackage[top=1in,bottom=1in,left=1in,right=1in]{geometry}
\usepackage{float}
\usepackage{graphicx}
\usepackage{mathtools}
\usepackage{tikz}
\usetikzlibrary{arrows}
\usepackage{circuitikz}

\date{
	\vspace{-.75em}
	{\large \today}
}

\title{
	\vspace*{-1.5em}
	{\Huge Lecture 6} \\
	\vspace{-1em}
	\rule{\textwidth/4}{.01em} \\
	\vspace{-.25em}
	{\Large Transfer Functions from Circuits}
	\vspace{-.75em}
}

\author{
	{\large Matthew Farah}
}

\hypersetup{
	colorlinks=true,
	linkcolor=blue,
	filecolor=magenta,
	urlcolor=blue,
	citecolor=blue,
	anchorcolor=blue
}

\fancyhead{}
\fancyhead[L]{\today}
\fancyhead[R]{Lecture 6 - Transfer Functions from Circuits}

\setlength{\parindent}{0cm}

\newcounter{example}
\renewcommand{\theexample}{\arabic{example}}

\newenvironment{example}[1][]{
	\refstepcounter{example}
	\begin{adjustwidth}{1.5em}{}
	\noindent\textbf{\ifx&#1&Example \theexample:\else Example \theexample \ (#1):\fi}
	\ignorespaces
}{
	\vspace{.5em}
	\end{adjustwidth}
}

\newenvironment{solution}{
	\vspace{.5em}
	\par
	\begin{adjustwidth}{1.5em}{}
	\textbf{{Solution:}}
	\par
	\vspace{.5em}
}{
	\vspace{1em}
	\end{adjustwidth}
}

\newcounter{subexample}
\renewcommand{\thesubexample}{\alph{subexample}}

\newenvironment{subexample}{
	\setcounter{step}{0}
	\refstepcounter{subexample}
	\vspace{1em}\par\noindent
	\begin{adjustwidth}{1.5em}{}
	\noindent\textbf{(\thesubexample)}
	\ignorespaces
}{
	\par
	\vspace{1em}
	\end{adjustwidth}
}

\renewenvironment{proof}{
	\vspace{.5em}\par\noindent
	\begin{adjustwidth}{1.5em}{}
	\emph{Proof:}\par\vspace{1em}
	\setlength{\parindent}{0cm}
}{
	\end{adjustwidth}
}

\newtheorem{theorem}{Theorem}

\renewenvironment{theorem}[1][]{
	\refstepcounter{theorem}
	\vspace{1em}\par\noindent
	\begin{adjustwidth}{1.5em}{}
	\noindent\textbf{\ifx&#1&Theorem \thetheorem:\else Theorem \thetheorem \ (#1):\fi}
	\ignorespaces
}{
	\vspace{.5em}
	\end{adjustwidth}
}

\newtheorem{lemma}{Lemma}

\renewenvironment{lemma}[1][]{
	\refstepcounter{lemma}
	\vspace{1em}\par\noindent
	\begin{adjustwidth}{1.5em}{}
	\noindent\textbf{\ifx&#1&Lemma \thelemma:\else Lemma \thelemma \ (#1):\fi}
	\ignorespaces
}{
	\vspace{.5em}
	\end{adjustwidth}
}

\makeatother
\makeatletter

\newcounter{step}
\renewcommand{\thestep}{\Roman{step}}

\newenvironment{step}{
	\refstepcounter{step}
	\vspace{1em}\par\noindent
	\ignorespaces\noindent\textbf{(\thestep)}
	\ignorespaces
}{
	\par
}

\sectionfont{\fontsize{12}{12}\bfseries}
\subsectionfont{\fontsize{10}{12}\normalfont}
\subsubsectionfont{\fontsize{10}{12}\normalfont}

\begin{document}

\maketitle
\pagestyle{fancy}

\fi

\begin{example}
	Find the transfer function for the voltage across the capacitor, $V_c(s)$, in the following circuit:
\end{example}

\vspace{1em}

\begin{figure}[H]
	\centering
	\begin{circuitikz}
		\draw (0,0) to[battery1, l=$V(s)$, invert] (0,3);
		\draw (0,3) to[L, l=$Ls$, i^>=$I(s)$] (3,3);
		\draw (3,3) to[R, l=$R$] (6,3);
		\draw (6,3) to[C, l_=$\tfrac{1}{Cs}$] (6,0);
		\node[right] at (6.5,1.5) {$V_c(s)$};
		\draw (6,0) -- (0,0);
	\end{circuitikz}
\end{figure}

\vspace{1em}

\begin{solution}
	We know that the transfer function representing the voltage across the capacitor is given by:

	\begin{align*}
		G_c(s) &= \frac{V_c(s)}{V(s)} \\
	\end{align*}

	By Kirchhoff's voltage law, we also know that the sum of the voltages around a loop must equal zero. So:

	\begin{align*}
		V(s) - LsI(s) - RI(s) - V_c(s) &= 0
	\end{align*}

	We also know the potential difference across the capacitor, $V_c(s)$, can be expressed by $V_c(s) = \frac{1}{Cs} I(s)$. So, using the above substitution and solving for $I(s)$, we find that:

	\begin{equation*}
		\begin{array}{lrl}
			& V(s) - LsI(s) - RI(s) - \tfrac{1}{Cs} I(s) =& 0 \\
			\implies & LsI(s) + RI(s) + \tfrac{1}{Cs} I(s) =& V(s) \\
			\implies & \left( Ls + R + \tfrac{1}{Cs} \right) I(s) =& V(s) \\
			\implies & I(s) =& \dfrac{V(s)}{Ls + R + \frac{1}{Cs}} \\
			\implies & I(s) =& \dfrac{Cs V(s)}{LCs^2 + RCs + 1} \\
		\end{array}
	\end{equation*}

	\vspace{0.5em}

	And since $V_c(s) = \frac{1}{Cs} I(s)$, as stated before, we can find $V_c(s)$ to be:

	\begin{align*}
		V_c(s) &= \tfrac{1}{Cs} I(s) \\
			&= \tfrac{1}{Cs} \left( \frac{Cs V(s)}{LCs^2 + RCs + 1} \right) \\
			&= \frac{V(s)}{LCs^2 + RCs + 1} \\
	\end{align*}

	And so, the transfer function for the voltage of the capacitor is given by:

	\begin{align*}
		G_c(s) &= \frac{V_c(s)}{V(s)} = \frac{1}{LCs^2 + RCs + 1} \\
	\end{align*}
\end{solution}

\begin{example}
	Find the transfer function, $G_c(s)$, for the voltage across the capacitor, $V_c(s)$, in the following circuit:
\end{example}

\vspace{1em}

\begin{figure}[H]
	\centering
	\begin{circuitikz}
		\draw (0,0) to[battery1, l=$V(s)$, invert] (0,3);
		\draw (0,3) to[R, l=$R_1$] (3,3);
		\draw (3,3) to[R, l=$R_2$] (6,3);
		\draw (3,3) to[L, l_=$Ls$] (3,0);
		\draw (6,3) to[C, l_=$\tfrac{1}{Cs}$] (6,0);
		\node[right] at (6.5,1.5) {$V_c(s)$};
		\draw (0,0) -- (3,0);
		\draw (3,0) -- (6,0);
		\draw[->, thick] (1.7,1.5) ++(135:0.5) arc (135:-135:0.5);
		\node at (1.7,1.5) {$I_1(s)$};
		\draw[->, thick] (4.3,1.5) ++(135:0.5) arc (135:-135:0.5);
		\node at (4.3,1.5) {$I_2(s)$};
	\end{circuitikz}
\end{figure}

\vspace{1em}

\begin{solution}
	As before, we know $G_c(s) = \frac{V_c(s)}{V(s)}$, and by inspection of the circuit, we can see that Kirchhoff's voltage law tells us that for the first loop:

	\begin{equation*}
		\begin{array}{lrl}
			& V(s) - R_1 I_1(s) - Ls(I_1(s) - I_2(s)) =& 0 \\
			\implies & R_1 I_1(s) + Ls(I_1(s) - I_2(s)) =& V(s) \\
			\implies & R_1 I_1(s) + Ls I_1(s) - Ls I_2(s) =& V(s) \\
			\implies & (R_1 + Ls) I_1(s) - Ls I_2(s) =& V(s) \quad \ldots (1) \\
		\end{array}
	\end{equation*}

	\vspace{0.5em}

	And similarly, for loop two:

	\begin{align*}
		Ls(I_2(s) - I_1(s)) + R_2 I_2(s) + V_c(s) &= 0 \\
		Ls I_2(s) - Ls I_1(s) + R_2 I_2(s) + V_c(s) &= 0 \\
	\end{align*}

	And since $V_c(s) = \frac{1}{Cs} I_2(s)$:

	\begin{equation*}
		\begin{array}{lrl}
			& Ls I_2(s) - Ls I_1(s) + R_2 I_2(s) + \tfrac{1}{Cs} I_2(s) =& 0 \\
			\implies & -Ls I_1(s) + \left( Ls + R_2 + \tfrac{1}{Cs} \right) I_2(s) =& 0 \quad \ldots (2) \\
		\end{array}
	\end{equation*}

	\vspace{0.5em}

	Putting equations $(1)$ and $(2)$ into matrix form, we can see that:

	\begin{align*}
		\begin{bmatrix} R_1 + Ls & -Ls \\ -Ls & Ls + R_2 + \tfrac{1}{Cs} \end{bmatrix} \begin{bmatrix} I_1(s) \\ I_2(s) \end{bmatrix} &= \begin{bmatrix} V(s) \\ 0 \end{bmatrix} \\
	\end{align*}

	Since we're interested in finding the transfer function for $V_c(s)$, and since $V_c(s) = \frac{1}{Cs} I_2(s)$, we should seek to compute $I_2(s)$. To do that, we will make use of Kramer's rule. For these particular matrices, Kramer's rule tells us that:

	\begin{align*}
		I_2(s) &= \frac{\det(A_2)}{\det(A)} \\
	\end{align*}

	Where:

	\begin{align*}
		A &= \begin{bmatrix} R_1 + Ls & -Ls \\ -Ls & Ls + R_2 + \tfrac{1}{Cs} \end{bmatrix}, \quad A_2 = \begin{bmatrix} R_1 + Ls & V(s) \\ -Ls & 0 \end{bmatrix} \\
	\end{align*}

	First, we can find that:

	\begin{align*}
		\det(A) &= (R_1 + Ls)\left( Ls + R_2 + \tfrac{1}{Cs} \right) - (-Ls)(-Ls) \\
			&= LR_1 s + R_1 R_2 + \tfrac{R_1}{Cs} + L^2 s^2 + L R_2 s + \tfrac{L}{C} - L^2 s^2 \\
			&= LR_1 s + R_1 R_2 + \tfrac{R_1}{Cs} + LR_2 s + \tfrac{L}{C} \\
			&= \tfrac{R_1}{Cs} + \left( R_1 R_2 + \tfrac{L}{C} \right) + (LR_1 + LR_2) s \\
	\end{align*}

	And similarly:

	\begin{align*}
		\det(A_2) &= (R_1 + Ls)(0) - (V(s))(-Ls) \\
			&= Ls V(s) \\
	\end{align*}

	Therefore:

	\begin{align*}
		I_2(s) &= \frac{Ls V(s)}{\frac{R_1}{Cs} + \left( R_1 R_2 + \frac{L}{C} \right) + (LR_1 + LR_2) s} \\
	\end{align*}

	Hence:

	\begin{align*}
		V_c(s) &= \tfrac{1}{Cs} I_2(s) \\
			&= \tfrac{1}{Cs} \left( \frac{Ls V(s)}{\frac{R_1}{Cs} + \left( R_1 R_2 + \frac{L}{C} \right) + (LR_1 + LR_2) s} \right) \\
			&= \frac{L V(s)}{R_1 + (R_1 R_2 C + L) s + (LR_1 C + LR_2 C) s^2} \\
	\end{align*}

	And thus:

	\begin{align*}
		G_c(s) &= \frac{V_c(s)}{V(s)} = \frac{L}{R_1 + (R_1 R_2 C + L) s + (LR_1 C + LR_2 C) s^2} \\
	\end{align*}
\end{solution}

\ifdefined\iscombined
	\newpage
\else
	\end{document}
\fi
